\section*{Две стенки в одном ГПУ-кристалле}
В ячейке \(12\times12\times8\) из 2304 узлов два независимых ансамбля имеют одинаковый MTP Qi, нулевое давление, свободные объём и \(c/a\). Меняется порог \(w\) для максимального теплового смещения. Для каждого атома \(i\) и каждого из двенадцати векторов идеальных соседей \(\mathbf d\) в текущей геометрии требуем
\[
 W=\max_{i,\mathbf d}\frac{2(\mathbf u_i-\bar{\mathbf u})\cdot\mathbf d}{|\mathbf d|^2}<w.
\]
Среднее смещение \(\bar{\mathbf u}\) исключает общий перенос кристалла. Стенка следует за масштабированием ячейки, вместе с ней меняются \(\mathbf d\) и физические смещения. При \(w=0{,}35\) отсекаются более крупные смещения, чем при \(w=0{,}40\). \href{https://doi.org/10.1103/PhysRevMaterials.7.103602}{Qi и соавт. (2023).}

\subsection*{Какие состояния исключает геометрическое условие}
В коде дополнительно заданы границы \(15\le v\le21\)~\AA$^3$/атом, \(1{,}45\le\rho=c/a\le1{,}72\) и проверка \(h_w=(1-w)d_{\rm nn}-1{,}68052~\text{\AA}\ge0\) для проб объёма и формы. Здесь \(d_{\rm nn}\) обозначает наименьшее расстояние между идеальными узлами,
\[
 a=\left(\frac{4v}{\sqrt3\,\rho}\right)^{1/3},\qquad
 d_{\rm nn}=a\min\!\left(1,\sqrt{1/3+\rho^2/4}\right).
\]
Наименьшее допустимое \(d_{\rm nn}=2{,}66928\)~\AA{} достигается при \(v=15\), \(\rho=1{,}45\). При \(w=0{,}35\) порог \(d_{\rm nn}\ge2{,}58542\)~\AA{} выполнен во всём прямоугольнике \((v,\rho)\). При \(w=0{,}40\) порог \(2{,}80087\)~\AA{} усекает область форм. Например, при \(\rho=1{,}60\) он поднимает нижнюю границу \(v\) с 15 до 15,853~\AA$^3$/атом. Поэтому формальное вложение ансамблей нужно проверять одновременно по координатам атомов и форме ячейки.

В каждом продуктивном интервале проверена нижняя граница \(h_{0{,}40}\) для всех возможных проб. Одна проба объёма изменяет \(v\) не более чем на \(0{,}22\sqrt{128/N}\)~\AA$^3$/атом, а одна проба формы изменяет \(\log\rho\) не более чем на \(0{,}008\sqrt{128/N}\) раз в три прохода. На фиксированном \(v\) наименьшее \(d_{\rm nn}\) внутри достижимого интервала \(\rho\) находится на одном из его концов. Если полученная нижняя граница положительна, геометрическое условие не могло отклонить ни одной геометрической пробы в интервале. Неположительная граница указывала бы только на возможность такого отказа.

\subsection*{Предыдущие наблюдения}
При 600~К сужение стенки \(0{,}40\to0{,}35\) на 1400 атомах изменило составы \(\alpha/\alpha_2\) с 5,866/24,265 до 5,872/24,313~ат.\% Al. При 1000~К те же составы изменились с 14,98/22,44 до 14,86/22,36. На 288 атомах при 1000~К более тесной стенке отвечали 15,8\% равновесных кадров \(\alpha\) и 46,8\% кадров \(\alpha_2\) в широком ансамбле. Вклад стенки в разность полубольших потенциалов \(\Omega_{\alpha_2}-\Omega_{\alpha}\) при исходном поле равен \(-k_{\rm B}T\ln(0{,}468/0{,}158)\), или \(-0{,}325\)~мэВ/атом.

При 1273,15~К, \(N=288\) и \(\Delta\mu=3{,}48\)~эВ сужение стенки повысило \(\langle Q\rangle\) с 0,29 до 0,40. Эти наблюдения дают основание ожидать увеличения веса упорядоченной фазы. В старой серии на 2304 атомах при 1245~К из 38\,400 продуктивных кадров при \(w=0{,}40\) ни один не попал в \(W<0{,}35\). Поэтому каждый новый порог исследован отдельными зонтичными траекториями.
